\section{Definisi Masalah}
\label{sec:masalah}

\subsection{Beban PLTU yang Tidak Ditanggung Pengelolanya}

PLTU Jawa-Bali terkonsentrasi di Banten, Jawa Tengah, dan Jawa Timur, sedangkan dampaknya menyebar ke kawasan padat penduduk (Tabel~\ref{tab:pltu}). Sebanyak 18,8 GW dimiliki IPP dan 12,0 GW dimiliki PLN. Dampak kesehatan tidak merata antar-PLTU. Untuk Cirebon-1, CREA memperkirakan pensiun pada 2035 menghindari 6.370 kematian selama 2036-2042, atau sekitar 910 kematian per tahun \cite{B-ECOBUSINESS-CIREBON}. Biaya kesehatan ini tidak masuk ke harga listrik.

\begin{table}[H]
\centering
\caption{PLTU terhubung jaringan PLN Jawa-Bali per provinsi}
\label{tab:pltu}
\small
\begin{tabular}{lrrrrrr}
\toprule
Provinsi & Unit & MW & IPP (MW) & PLN (MW) & Mt CO$_2$/tahun & Kematian/tahun \\
\midrule
Banten & 19 & 10.527 & 4.642 & 5.885 & 57,6 & 2.723 \\
Jawa Tengah & 15 & 10.090 & 8.800 & 1.290 & 52,4 & 2.610 \\
Jawa Timur & 12 & 6.155 & 3.365 & 2.790 & 36,3 & 1.592 \\
Jawa Barat & 8 & 3.624 & 1.584 & 2.040 & 19,9 & 1.676 \\
Bali & 3 & 426 & 426 & 0 & 2,7 & 110 \\
\midrule
Total & 57 & 30.822 & 18.817 & 12.005 & 168,9 & 8.711 \\
\bottomrule
\end{tabular}

\sumber{Global Energy Monitor \cite{D-GEM-GCPT}. Kematian dini dari faktor rata-rata armada CREA \cite{P-CREA-IESR-2023}, kecuali Cirebon-1 \cite{B-ECOBUSINESS-CIREBON}}
\end{table}

\subsection{Kontrak dan Kompensasi Menghambat Pensiun Dini}

PLN membayar pembayaran kapasitas kepada IPP selama unit tersedia, terlepas dari listriknya dipakai atau tidak \cite{P-IEEFA-2017}. Mengakhiri kontrak lebih awal menimbulkan kewajiban kompensasi atas keuntungan yang hilang. Pensiun sebelum aset habis disusutkan juga menimbulkan risiko aset terdampar \cite{P-KEFFORD-2018}. Pada kasus Cirebon-1, ADB menyiapkan refinancing US\$250-300 juta untuk percepatan sekitar tujuh tahun \cite{B-REUTERS-CIREBON-2022,B-ADB-CIREBON}, sedangkan PLN menilai kewajibannya mencapai sekitar Rp60 triliun dan memilih pendekatan phase down \cite{B-PETROMINDO-CIREBON-BATAL,B-MONGABAY-202601}. Selisih penilaian sebesar ini menunjukkan bahwa besaran kompensasi adalah titik perdebatan utama.

\subsection{PLTS Atap Terhambat Aturan dan Modal Awal}

Permen ESDM 2/2024 menghapus perhitungan ekspor listrik sebagai pengurang tagihan dan menetapkan kuota per wilayah \cite{R-PERMEN-ESDM-2-2024}. Tanpa nilai ekspor, PLTS atap rumah tangga kurang menarik karena konsumsi rumah tangga tinggi pada malam hari, sehingga adopsi bergeser ke sektor komersial dan industri \cite{P-IESR-PERMEN2,B-PVMAG-QUOTA}. Harga pemasangan masih berada pada kisaran belasan juta rupiah per kWp \cite{B-DETIK-HARGA-PLTS}. Realisasi 538 MWp pada Juli 2025 masih jauh dari kuota 2 GW pada 2028 \cite{B-KONTAN-ROOFTOP-2025,B-CNBC-ROOFTOP-2028}.

\subsection{Beban Fiskal Berpindah ke APBN}

Tarif listrik ditetapkan pemerintah di bawah biaya pokok penyediaan (BPP). Subsidi listrik 2026 diusulkan Rp97,37-104,97 triliun untuk 44,88 juta pelanggan \cite{B-KONTAN-SUBSIDI-2026}. Setiap tambahan biaya PLN pada akhirnya berpindah ke APBN. Ketika PLTS atap mengurangi penjualan PLN, pendapatan PLN turun, sedangkan pembayaran kapasitas PLTU tetap berjalan. Kondisi inilah yang membuat PLN menolak nilai ekspor PLTS atap selama sistem Jawa-Bali kelebihan pasokan \cite{B-TENGGARA-2024}.

\subsection{Pekerja dan Daerah Terdampak}

Cirebon-1 mempekerjakan sekitar 300 orang. CELIOS memperkirakan sedikitnya 14.000 pekerja langsung dan tidak langsung terdampak bila Cirebon-1 dan Pelabuhan Ratu dipensiunkan \cite{B-ECOBUSINESS-CIREBON}. Sebaliknya, PLTS atap termasuk teknologi energi paling padat karya, sekitar 24,72 job-year per MW \cite{P-CEEW-2017}. Peluang ini tidak otomatis menyerap pekerja PLTU tanpa program pelatihan ulang yang dijadwalkan.
